% GENERATED FILE. Edit canonical lessons, then run npm run generate:books.
% canonical-source-hash: fnv1a64:960d6acb5de01fd1
% canonical-lessons: PA-C272-vichharna, PA-C272-sipahi, PA-C272-chauki, PA-C272-riksha, PA-C272-parak

\chapter{Lost in the Crowd}
\label{ch:pa-a2-lost-in-the-crowd}
\hlchaptermodality{272}

\begin{chapteropening}
\textbf{By the end of this chapter:} \emph{I can ask for help when I lose my group in town.}

Complete the last lesson of chapter 272: I can ask for help when I lose my group in town.
\end{chapteropening}

\section[\texorpdfstring{\pa{ਵਿਛੜਨਾ} (vichhaṛnā)}{ਵਿਛੜਨਾ (vichhaṛnā)}]{\pa{ਵਿਛੜਨਾ} (vichhaṛnā) — to get separated; to go astray}

\label{lesson:PA-C272-vichharna}



\begin{quote}
Before the new one: say the Punjabi for west, then the Punjabi for to show someone around.
\end{quote}

\subsection*{You'll want to know: \pa{ਵਿਛੜਨਾ}}
\textbf{\pa{ਵਿਛੜਨਾ}} — \emph{vichhaṛnā} — \textquotedblleft{}to get separated; to go astray\textquotedblright{}.

\begin{quote}
\pa{ਭੀੜ} \pa{ਵਿੱਚ} \pa{ਸਾਡੇ} \pa{ਤੋਂ} \pa{ਨਾ} \pa{ਵਿਛੜੋ।}
\end{quote}

\begin{quote}
\emph{bhīṛ vich sāḍe ton nā vichhṛo.}
\end{quote}

\begin{quote}
Do not get separated from us in the crowd.
\end{quote}

Past (I …): \pa{ਮੈਂ} \pa{ਦੋਸਤਾਂ} \pa{ਤੋਂ} \pa{ਵਿਛੜਿਆ}, \emph{maiṁ dostāṁ ton vichhṛiā}.

\begin{cousinweb}
\textbf{\pa{ਵਿਛੜਨਾ}} (\emph{vichhaṛnā}) is inherited from a reconstructed \emph{*vicchaṭ}, \textquotedblleft{}be separated\textquotedblright{}; CDIAL lists Punjabi \emph{vicchaṛṇā}, \textquotedblleft{}to be separated, go astray\textquotedblright{}, under it (CDIAL 11687).
\end{cousinweb}

\begin{grammarlens}[title={What you've built}]
One more word for this chapter.
\end{grammarlens}

\subsection*{Your turn}
\begin{itemize}
\raggedright
  \item \emph{Say it:} \emph{vichhaṛnā}
  \item \emph{Say it:} \emph{vichhaṛnā}, once more
  \item \emph{Say it:} \emph{bhīṛ vich sāḍe ton nā vichhṛo.} — Do not get separated from us in the crowd.
  \item \emph{Recall:} say the Punjabi for west, then the Punjabi for to show someone around, then say \emph{vichhaṛnā} again
\end{itemize}

\begin{tcolorbox}[breakable,colback=teal!4,colframe=teal!35!black,title={Before you move on}]
What does \pa{ਵਿਛੜਨਾ} mean? (\textquotedblleft{}To get separated; to go astray\textquotedblright{}.) Say it once more.
\end{tcolorbox}
\section[\texorpdfstring{\pa{ਸਿਪਾਹੀ} (sipāhī)}{ਸਿਪਾਹੀ (sipāhī)}]{\pa{ਸਿਪਾਹੀ} (sipāhī) — a police constable; a soldier}

\label{lesson:PA-C272-sipahi}



\begin{quote}
Before the new one: say the Punjabi for to show someone around, then the Punjabi for to get separated.
\end{quote}

\subsection*{You'll want to know: \pa{ਸਿਪਾਹੀ}}
\textbf{\pa{ਸਿਪਾਹੀ}} — \emph{sipāhī} — \textquotedblleft{}a police constable; a soldier\textquotedblright{}.

\begin{quote}
\pa{ਚੌਂਕ} \pa{ਵਿੱਚ} \pa{ਇੱਕ} \pa{ਸਿਪਾਹੀ} \pa{ਹੈ।}
\end{quote}

\begin{quote}
\emph{chauṅk vich ikk sipāhī hai.}
\end{quote}

\begin{quote}
There is a constable in the square.
\end{quote}

A masculine noun.

\begin{cousinweb}
\textbf{\pa{ਸਿਪਾਹੀ}} (\emph{sipāhī}) is a loan, usually traced to Persian \emph{sipāhī}, \textquotedblleft{}a soldier\textquotedblright{}.
\end{cousinweb}

\begin{grammarlens}[title={What you've built}]
One more word for this chapter.
\end{grammarlens}

\subsection*{Your turn}
\begin{itemize}
\raggedright
  \item \emph{Say it:} \emph{sipāhī}
  \item \emph{Say it:} \emph{sipāhī}, once more
  \item \emph{Say it:} \emph{chauṅk vich ikk sipāhī hai.} — There is a constable in the square.
  \item \emph{Recall:} say the Punjabi for to show someone around, then the Punjabi for to get separated, then say \emph{sipāhī} again
  \item \emph{Write it:} copy \textbf{\pa{ਸਿਪਾਹੀ}} once, checking the whole word
\end{itemize}

\begin{tcolorbox}[breakable,colback=teal!4,colframe=teal!35!black,title={Before you move on}]
What does \pa{ਸਿਪਾਹੀ} mean? (\textquotedblleft{}A police constable; a soldier\textquotedblright{}.) Say it once more.
\end{tcolorbox}
\section[\texorpdfstring{\pa{ਚੌਕੀ} (chaukī)}{ਚੌਕੀ (chaukī)}]{\pa{ਚੌਕੀ} (chaukī) — a police post; a guard post}

\label{lesson:PA-C272-chauki}



\begin{quote}
Before the new one: say the Punjabi for to get separated, then the Punjabi for a police constable.
\end{quote}

\subsection*{You'll want to know: \pa{ਚੌਕੀ}}
\textbf{\pa{ਚੌਕੀ}} — \emph{chaukī} — \textquotedblleft{}a police post; a guard post\textquotedblright{}.

\begin{quote}
\pa{ਪੁਲਿਸ} \pa{ਦੀ} \pa{ਚੌਕੀ} \pa{ਪੁਲ} \pa{ਦੇ} \pa{ਨੇੜੇ} \pa{ਹੈ।}
\end{quote}

\begin{quote}
\emph{pulis dī chaukī pul de neṛe hai.}
\end{quote}

\begin{quote}
The police post is near the bridge.
\end{quote}

A feminine noun.

\begin{cousinweb}
\textbf{\pa{ਚੌਕੀ}} (\emph{chaukī}) is inherited from Sanskrit \emph{catuṣka}, \textquotedblleft{}a set of four\textquotedblright{}; CDIAL lists Punjabi \emph{caukī}, \textquotedblleft{}chair\textquotedblright{}, under it (CDIAL 4629). The sense \textquotedblleft{}a police post\textquotedblright{} is not given there.
\end{cousinweb}

\begin{grammarlens}[title={What you've built}]
One more word for this chapter.
\end{grammarlens}

\subsection*{Your turn}
\begin{itemize}
\raggedright
  \item \emph{Say it:} \emph{chaukī}
  \item \emph{Say it:} \emph{chaukī}, once more
  \item \emph{Say it:} \emph{pulis dī chaukī pul de neṛe hai.} — The police post is near the bridge.
  \item \emph{Recall:} say the Punjabi for to get separated, then the Punjabi for a police constable, then say \emph{chaukī} again
  \item \emph{Write it:} copy \textbf{\pa{ਚੌਕੀ}} once, checking the whole word
\end{itemize}

\begin{tcolorbox}[breakable,colback=teal!4,colframe=teal!35!black,title={Before you move on}]
What does \pa{ਚੌਕੀ} mean? (\textquotedblleft{}A police post; a guard post\textquotedblright{}.) Say it once more.
\end{tcolorbox}
\section[\texorpdfstring{\pa{ਰਿਕਸ਼ਾ} (rikshā)}{ਰਿਕਸ਼ਾ (rikshā)}]{\pa{ਰਿਕਸ਼ਾ} (rikshā) — a rickshaw}

\label{lesson:PA-C272-riksha}



\begin{quote}
Before the new one: say the Punjabi for a police constable, then the Punjabi for a police post.
\end{quote}

\subsection*{You'll want to know: \pa{ਰਿਕਸ਼ਾ}}
\textbf{\pa{ਰਿਕਸ਼ਾ}} — \emph{rikshā} — \textquotedblleft{}a rickshaw\textquotedblright{}.

\begin{quote}
\pa{ਸਟੇਸ਼ਨ} \pa{ਦੇ} \pa{ਬਾਹਰ} \pa{ਬਹੁਤ} \pa{ਰਿਕਸ਼ੇ} \pa{ਹਨ।}
\end{quote}

\begin{quote}
\emph{saṭeshan de bāhar bahut rikshe han.}
\end{quote}

\begin{quote}
There are many rickshaws outside the station.
\end{quote}

A masculine noun.

\begin{cousinweb}
\textbf{\pa{ਰਿਕਸ਼ਾ}} (\emph{rikshā}) is borrowed from English \emph{rickshaw}.
\end{cousinweb}

\begin{grammarlens}[title={What you've built}]
One more word for this chapter.
\end{grammarlens}

\subsection*{Your turn}
\begin{itemize}
\raggedright
  \item \emph{Say it:} \emph{rikshā}
  \item \emph{Say it:} \emph{rikshā}, once more
  \item \emph{Say it:} \emph{saṭeshan de bāhar bahut rikshe han.} — There are many rickshaws outside the station.
  \item \emph{Recall:} say the Punjabi for a police constable, then the Punjabi for a police post, then say \emph{rikshā} again
  \item \emph{Write it:} copy \textbf{\pa{ਰਿਕਸ਼ਾ}} once, checking the whole word
\end{itemize}

\begin{tcolorbox}[breakable,colback=teal!4,colframe=teal!35!black,title={Before you move on}]
What does \pa{ਰਿਕਸ਼ਾ} mean? (\textquotedblleft{}A rickshaw\textquotedblright{}.) Say it once more.
\end{tcolorbox}
\section[\texorpdfstring{\pa{ਪਾਰਕ} (pārak)}{ਪਾਰਕ (pārak)}]{\pa{ਪਾਰਕ} (pārak) — a park}

\label{lesson:PA-C272-parak}



\begin{quote}
Before the new one: say the Punjabi for a police post, then the Punjabi for a rickshaw.
\end{quote}

\subsection*{You'll want to know: \pa{ਪਾਰਕ}}
\textbf{\pa{ਪਾਰਕ}} — \emph{pārak} — \textquotedblleft{}a park\textquotedblright{}.

\begin{quote}
\pa{ਬੱਚੇ} \pa{ਸ਼ਾਮ} \pa{ਨੂੰ} \pa{ਪਾਰਕ} \pa{ਵਿੱਚ} \pa{ਖੇਡਦੇ} \pa{ਹਨ।}
\end{quote}

\begin{quote}
\emph{bachche shām nū̃ pārak vich kheḍde han.}
\end{quote}

\begin{quote}
The children play in the park in the evening.
\end{quote}

A masculine noun.

\begin{cousinweb}
\textbf{\pa{ਪਾਰਕ}} (\emph{pārak}) is borrowed from English \emph{park}.
\end{cousinweb}

\begin{grammarlens}[title={What you've built}]
One more word for this chapter.
\end{grammarlens}

\subsection*{Your turn}
\begin{itemize}
\raggedright
  \item \emph{Say it:} \emph{pārak}
  \item \emph{Say it:} \emph{pārak}, once more
  \item \emph{Say it:} \emph{bachche shām nū̃ pārak vich kheḍde han.} — The children play in the park in the evening.
  \item \emph{Recall:} say the Punjabi for a police post, then the Punjabi for a rickshaw, then say \emph{pārak} again
  \item \emph{Write it:} copy \textbf{\pa{ਪਾਰਕ}} once, checking the whole word
\end{itemize}

\begin{tcolorbox}[breakable,colback=teal!4,colframe=teal!35!black,title={Before you move on}]
What does \pa{ਪਾਰਕ} mean? (\textquotedblleft{}A park\textquotedblright{}.) Say it once more.
\end{tcolorbox}
